\documentclass[a4paper,11pt]{article}
\usepackage[utf8]{inputenc}
\usepackage[T1]{fontenc}
\usepackage[ngerman]{babel}
\usepackage{amsmath,amssymb}
\usepackage[margin=2.3cm]{geometry}
\usepackage{enumitem}
\usepackage{fancyhdr}
\usepackage{array}
\usepackage[table]{xcolor}
\usepackage{tikz}
\definecolor{bgdark}{HTML}{1E1E1E}
\definecolor{fgtext}{HTML}{E6E6E6}
\definecolor{gridcol}{HTML}{4A4A4A}
\pagecolor{bgdark}
\arrayrulecolor{fgtext}
\AtBeginDocument{\color{fgtext}}
\renewcommand{\headrule}{\color{fgtext}\hrule height \headrulewidth}
\newcommand{\karo}[1]{\par\noindent
  \begin{tikzpicture}
    \draw[step=5mm, gridcol, very thin] (0,0) grid (\linewidth,#1);
  \end{tikzpicture}\par}

\newcommand{\diff}[1]{\hfill\normalfont\textbf{[#1]}}
\newcommand{\aufg}[2]{\subsection*{Aufgabe #1 \diff{#2}}}
\renewcommand{\arraystretch}{1.2}

\pagestyle{fancy}
\fancyhf{}
\lhead{RADM -- Mengen, Relationen, Abbildungen}
\rhead{Übungsblatt 02}
\cfoot{\thepage}

\begin{document}

\begin{center}
  {\LARGE\bfseries Übungsaufgaben Mengen, Relationen, Abbildungen}\\[2pt]
  {\large Relationen, Algebra, Diskrete Mathematik -- Teil 3}\\[10pt]
\end{center}

\noindent
\begin{tabular}{@{}p{0.6\textwidth}@{}p{0.4\textwidth}@{}}
  \textbf{Name:} \hrulefill & \textbf{Datum:} \hrulefill
\end{tabular}

\vspace{4pt}
\noindent\rule{\textwidth}{0.4pt}
\vspace{4pt}

\noindent\textit{Themen: Mengenoperationen, Mächtigkeit/Potenzmenge, DeMorgan,
kartesisches Produkt, algebraische Strukturen (Gruppe), Relationen \& ihre
Eigenschaften, Äquivalenz-/Ordnungsrelationen, Abbildungen, Mengenbeweise,
Relationenalgebra/Datenbanken.}

\vspace{6pt}

% ---------------------------------------------------------------
\aufg{1}{leicht}
Grundmenge $E=\{1,2,\dots,9\}$, sowie $A=\{1,2,3,4,5,6\}$, $B=\{2,4,6,8\}$, $C=\{1,3,5\}$.
Bestimme:
\[ A\cap B,\quad A\cup C,\quad A\setminus B,\quad B\setminus A,\quad \overline{C}\ (=E\setminus C),\quad (A\cap B)\cup C. \]
\karo{5cm}

% ---------------------------------------------------------------
\aufg{2}{leicht}
Sei $M=\{a,b,c\}$.
\begin{enumerate}[label=(\alph*)]
  \item Gib die Potenzmenge $\mathcal{P}(M)$ vollständig an und bestimme $|\mathcal{P}(M)|$.
  \item Bestimme $|\mathcal{P}(\mathcal{P}(\emptyset))|$.
\end{enumerate}
\karo{5cm}

% ---------------------------------------------------------------
\aufg{3}{mittel}
$E=\{1,\dots,10\}$, $A=\{1,2,3,4,5\}$, $B=\{4,5,6,7\}$. Verifiziere die DeMorgansche
Regel $\overline{A\cup B}=\overline{A}\cap\overline{B}$ durch Berechnung \emph{beider} Seiten.
\karo{5.5cm}

% ---------------------------------------------------------------
\aufg{4}{mittel}
$M=\{a,b\}$, $N=\{1,2,3\}$.
\begin{enumerate}[label=(\alph*)]
  \item Gib $M\times N$ an und bestimme $|M\times N|$.
  \item Gib $N\times M$ an. Gilt $M\times N = N\times M$?
\end{enumerate}
\karo{5.5cm}

% ---------------------------------------------------------------
\aufg{5}{mittel}
Untersuche, ob $(\mathbb{Z},-)$ -- die ganzen Zahlen mit der Subtraktion -- eine Gruppe ist.
Prüfe die Gruppenaxiome (Abgeschlossenheit, Assoziativität, neutrales Element, inverse Elemente).
\karo{6cm}

% ---------------------------------------------------------------
\aufg{6}{mittel}
$M=\{1,2,3,4\}$ und $R=\{(x,y)\in M^2 \mid x \text{ teilt } y \text{ ganzzahlig}\}$.
\begin{enumerate}[label=(\alph*)]
  \item Liste alle Paare von $R$ auf.
  \item Prüfe: reflexiv, symmetrisch, antisymmetrisch, transitiv.
  \item Um welche Art von Relation handelt es sich?
\end{enumerate}
\karo{6cm}

% ---------------------------------------------------------------
\aufg{7}{mittel}
Auf $\mathbb{Z}$ sei $x\sim y :\Leftrightarrow x-y$ ist durch $3$ teilbar.
\begin{enumerate}[label=(\alph*)]
  \item Zeige, dass $\sim$ eine Äquivalenzrelation ist.
  \item Gib alle Äquivalenzklassen an.
\end{enumerate}
\karo{6cm}

% ---------------------------------------------------------------
\aufg{8}{mittel}
Untersuche jede Abbildung auf \emph{injektiv}, \emph{surjektiv}, \emph{bijektiv}
(beachte Definitions- und Wertemenge!). Hier $\mathbb{N}=\{1,2,3,\dots\}$.
\begin{enumerate}[label=(\alph*)]
  \item $f:\mathbb{R}\to\mathbb{R}$, $f(x)=2x+1$.
  \item $f:\mathbb{R}\to\mathbb{R}$, $f(x)=x^2$.
  \item $f:\mathbb{R}\to\mathbb{R}_{\ge 0}$, $f(x)=x^2$.
  \item $f:\mathbb{N}\to\mathbb{N}$, $f(x)=x+1$.
\end{enumerate}
\karo{6cm}

% ---------------------------------------------------------------
\aufg{9}{schwer}
Beweise für beliebige Mengen $A,B,C$:
\[ A\setminus(B\cup C) = (A\setminus B)\cap(A\setminus C). \]
(Hinweis: Element-Beweis über die zugehörigen aussagenlogischen Bedingungen.)
\karo{6.5cm}

% ---------------------------------------------------------------
\aufg{10}{mittel}
Gegeben zwei Relationen (Tabellen):
\[
\textbf{Mitglieder}(\text{Nr},\text{Name}) = \{(1,\text{Anna}),(2,\text{Bernd}),(3,\text{Carla})\}
\]
\[
\textbf{Sportart}(\text{Nr},\text{Sport}) = \{(1,\text{Fußball}),(1,\text{Turnen}),(2,\text{Turnen}),(3,\text{Tennis})\}
\]
Bestimme:
\begin{enumerate}[label=(\alph*)]
  \item $\sigma_{\text{Sport}=\text{Turnen}}(\textbf{Sportart})$
  \item $\Pi_{\text{Name}}(\textbf{Mitglieder})$
  \item $\Pi_{\text{Name},\text{Sport}}(\textbf{Mitglieder}\bowtie\textbf{Sportart})$ \quad(natural Join über Nr)
\end{enumerate}
\karo{6cm}

% ---------------------------------------------------------------
\aufg{11}{leicht}
$M=\{1,2,3,4,5\}$. Bestimme den Wahrheitswert:
\begin{enumerate}[label=(\alph*)]
  \item $\forall x\in M:\; x\le 5$
  \item $\exists x\in M:\; x>4$
  \item $\forall x\in M:\; x \text{ ist gerade}$
  \item $\exists x\in M:\; x^2=9$
\end{enumerate}
\karo{4.5cm}

\end{document}
